\documentclass[11pt]{article}
\usepackage{mystyle}

\title{Rosen --- \S 4.1: Divisibility and Modular Arithmetic\\ \large Selected Exercises}
\author{Alston}
\date{Semester 2, 2026}

\begin{document}
\maketitle

% ======================================================================
\begin{exercise}{5}
    Show that if $a \mid b$ and $b \mid a$, where $a$ and $b$ are
    integers, then $a = b$ or $a = -b$.
\end{exercise}

% ======================================================================
\begin{exercise}{7}
    Show that if $a$, $b$, and $c$ are integers, where $a \neq 0$ and
    $c \neq 0$, such that $ac \mid bc$, then $a \mid b$.
\end{exercise}

% ======================================================================
\begin{exercise}{8}
    Prove or disprove that if $a \mid bc$, where $a$, $b$, and $c$
    are positive integers and $a \neq 0$, then $a \mid b$ or
    $a \mid c$.
\end{exercise}

% ======================================================================
\begin{exercise}{17}
    Show that if $n$ and $k$ are positive integers, then
    $\lceil n/k \rceil = \lfloor (n-1)/k \rfloor + 1$.
\end{exercise}

% ======================================================================
\begin{exercise}{19}
    Find a formula for the integer with smallest absolute value that
    is congruent to an integer $a$ modulo $m$, where $m$ is a
    positive integer.
\end{exercise}

\end{document}